\section{Polynomial likelihood reduction of a bounded-degree parent}
\label{plc:section}

Let $Y_t$ be the retained descendant component of a newborn tag, viewed
with its parent edge replaced by an inert stub.  Let $D_t$ be the full
degree of its real external parent.  Fix $M\ge2$ and an initial degree
$d\in\{2,\ldots,M\}$, and stop both coordinates at
\[
 \tau_M=\inf\{t:D_t=M+1\}.
\]
Before this stopping time, the parent's degree-increase intensity is
\[
 \lambda_p(t)=\sum_{v\sim p}\frac1{D_v(t)}\le D_t\le M.
\]
The sum includes real leaf neighbours, which contribute one, and
inert neighbours, which contribute zero.  Since the tag is a real
neighbour, $\lambda_p(t)>0$ at every finite state.
There are at most $K=M+1-d$ parent degree increases before stopping.

\begin{theorem}[A finite family of reference laws with polynomial loss]
\label{plc:theorem}
There is a reference law $\mathbb Q_d$ for $(Y,D)$, depending only on
$d$ and $M$, with time-dependent generator before stopping
\begin{align}
 \mathcal G_t^{\rm ref}f(y,d)
 &=\mathcal L_\infty f(\cdot,d)(y)
   +\frac1d\{f(y^+,d)-f(y,d)\}\notag\\
 &\quad+\frac1{(1+t)^2}\{f(y,d+1)-f(y,d)\},
 \label{plc:reference-generator}
\end{align}
such that every actual finite exterior with parent degree $d$ satisfies
the following bound.  For every $T\ge0$ and every nonnegative measurable
functional $Z$ of the stopped retained-and-degree path through time $T$,
\begin{equation}\label{plc:domination}
 \mathbb E_{
m actual} Z
 \le e\,[M(1+T)^2]^K\,\mathbb E_{\mathbb Q_d} Z.
\end{equation}
Thus arbitrary exterior histories with a bounded-degree parent can be
replaced by one of finitely many explicit reference models at a
polynomial, rather than exponential, cost in the time horizon.
\end{theorem}
\begin{proof}
Use the full finite-tree birth process until $\tau_M$.  A parent degree
increase is exactly one leaf addition at the parent, with total rate
$\lambda_p(t)$.  Change this single transition channel to the rate
$h(t)=(1+t)^{-2}$ and leave every other transition rate unchanged.
Both the original and the modified processes are nonexplosive: the
modified total real birth rate is at most the original population plus
one.  The usual finite-horizon counting-process change of measure gives
the actual law relative to the modified law with density
\[
 Z_T^{\rm RN}
 =\exp\!\left\{\int_0^{T\wedge\tau_M}
                     (h(t)-\lambda_p(t))\,dt\right\}
   \prod_{j:\,\sigma_j\le T\wedge\tau_M}
        \frac{\lambda_p(\sigma_j-)}{h(\sigma_j)}.
\]
Here $\sigma_j$ are the parent degree-increase times, including the
stopping increase when it occurs before $T$.  Localization in total
population first reduces the calculation to bounded finite-state
intensities.  Alternatively, the density follows directly by taking
the ratio of the finite jump-history densities and then passing to the
nonexplosive limit.

There are at most $K$ factors in the product,
$\lambda_p\le M$, and $\int_0^\infty h(t)dt=1$.  Hence pathwise
\[
 Z_T^{\rm RN}\le e\,[M(1+T)^2]^K.
\]

It remains to identify the marginal law of the retained coordinates
under the modified full process.  A parent degree increase changes only
$D$, and now has deterministic hazard $h(t)$.  A parent selection of the
tag creates one retained root leaf at rate $1/D$.  All other retained
transitions have exactly the inert generator $\mathcal L_\infty$.
Every remaining exterior transition changes neither $Y$ nor $D$.
Thus the pair $(Y,D)$ has precisely the autonomous generator
\eqref{plc:reference-generator}; in particular its law no longer depends
on the rest of the initial exterior.  This proves
\eqref{plc:domination}.
\end{proof}

The reference degree chain is driven by an independent inhomogeneous
Poisson clock with integrated hazard one, stopped after $K$ jumps.
The input $1/D_t$ is nonincreasing, changes at most $K$ times, and is
eventually constant on every path that does not reach the stopping
state.  For example, taking
$Z=C(Y_T)\mathbf1_{\{T<\tau_M\}}$ gives a uniform bound for the part
of the tagged mean that remains behind a low-degree parent.

More generally one may use
$h(t)=\eta(1+t)^{-\beta}$ with $\eta>0$ and $\beta>1$.
The likelihood factor is then at most
\[
 \exp\!\left(\frac{\eta}{\beta-1}\right)
 \max\left\{1,\frac M\eta(1+T)^\beta\right\}^K.
\]
Thus the total reference jump intensity can be made arbitrarily small
without introducing an exponential loss in~$T$.

\subsection{An exponential occupation bound after the root is large}

\begin{proposition}[Uniform inverse-degree exponential moment]
\label{plc:occupation}
Suppose a retained root starts with full degree $D_0\ge H\ge2$ and
receives incoming leaves at predictable intensity at least
$\rho\in(0,1]$.  The other incoming events may be arbitrary events of
the usual core process.  Put
\[
 c_\rho=1+\frac3\rho,\qquad
 \kappa_\rho=\frac\rho3,\qquad
 A_{H,\rho}=\frac{\log(c_\rho H)+1}{\kappa_\rho H}.
\]
Then
\begin{equation}\label{plc:occupation-mean}
 \mathbb E\int_0^\infty\frac{dt}{D_t}\le A_{H,\rho}.
\end{equation}
For every $p>0$ with $pA_{H,\rho}<1$,
\begin{equation}\label{plc:occupation-exponential}
 \boxed{
 \mathbb E\exp\!\left(p\int_0^\infty\frac{dt}{D_t}\right)
 \le\frac1{1-pA_{H,\rho}}.
 }
\end{equation}
The estimates are uniform over all finite internal environments, initial
leaf counts and other permitted inputs.  They also hold with the
integrals stopped at any stopping time up to which the lower immigration
bound remains valid.
\end{proposition}
\begin{proof}
The root degree is nondecreasing, so $D_t\ge H$.  Lower-couple its
leaf count with the birth--death chain $X$ started at zero, with births
at rate $x+\rho$ and deaths at rate one whenever $x>0$.  A rate-$\rho$
stream is extracted from the incoming leaves, the leaf births are
coupled by their Poisson random measures, and the comparison loses
one leaf at every root clock ring.  Actual root leaf losses occur at
most once at such a ring.  This is the same comparison used in
Lemma~\ref{ps:reciprocal}.  Its reciprocal estimate gives
\[
 \mathbb E\frac1{D_t}
 \le\mathbb E\min\left\{\frac1H,\frac1{X_t+1}\right\}
 \le\min\left\{\frac1H,c_\rho e^{-\kappa_\rho t}\right\}.
\]
Integrating the last bound, with the two terms equal at
$t=\log(c_\rho H)/\kappa_\rho$, proves
\eqref{plc:occupation-mean}.

The same bound applies conditionally after every stopping time: the
root degree is still at least $H$, the comparison chain may be restarted
at zero below its actual leaf count, and the remaining input still has
intensity at least $\rho$.  For
$I=\int_0^\infty D_t^{-1}dt$, iterated conditioning in the ordered
time integrals therefore gives the Khasminskii bounds
\[
 \mathbb E I^n\le n!A_{H,\rho}^n\qquad(n\ge0).
\]
Summing the nonnegative exponential series proves
\eqref{plc:occupation-exponential}.  With a stopping time, insert the
indicator that the input lower bound is still in force; the same
conditional argument only decreases the integrals.
\end{proof}

In particular, $A_{H,\rho}\to0$ as $H\to\infty$ for fixed $\rho$.
Thus any prescribed exponential moment is uniformly finite once the
root degree is sufficiently large.

